\documentclass[12pt, letterpaper]{article}

% --- PAQUETES ---
\usepackage[utf8]{inputenc}
\usepackage[T1]{fontenc}
\usepackage[spanish]{babel}
\usepackage[left=2cm, right=2cm, top=2cm, bottom=2.5cm]{geometry}
\usepackage{enumitem}
\usepackage{amssymb}
\usepackage{tcolorbox}
\usepackage{amsmath}
\usepackage{ifthen}
\usepackage{xcolor}

% --- FUENTE --- 
\usepackage{helvet} 
\renewcommand{\familydefault}{\sfdefault}

% ==========================================
%  CONFIGURACIÓN DE RESPUESTAS (ACTIVAR/DESACTIVAR)
% ==========================================
% Cambia 'true' por 'false' para ocultar las respuestas e imprimir la versión del estudiante.
\newboolean{mostrarRespuestas}
\setboolean{mostrarRespuestas}{true}

% Comando interno para generar las cajas de solución
\newcommand{\solucion}[1]{
	\ifthenelse{\boolean{mostrarRespuestas}}{
		\vspace{0.2cm}
		\begin{tcolorbox}[colframe=red!70!black, colback=red!5, boxrule=1pt, arc=3pt, left=5pt, right=5pt, top=5pt, bottom=5pt]
			\textcolor{red!70!black}{\textbf{Solución de calificación:} #1}
		\end{tcolorbox}
	}{}
}

\begin{document}
	
	% --- ENCABEZADO --- 
	\begin{center} 
		{\Large \textbf{LICEO V.A.L. (Vida, Amor, Luz)}}\\[0.2cm] 
		{\large \textbf{DEPARTAMENTO DE CIENCIAS BÁSICAS}}\\[0.2cm] 
		\textbf{Taller de Repaso - CIENCIAS 2}\\[0.1cm] 
		\textit{TAU 3: Entre la Luna y la Tierra} 
	\end{center}
	
	\vspace{0.3cm}
	\noindent \textbf{Nombre:} \underline{\hspace{8cm}} \hspace{1cm}\textbf{Fecha:} \underline{\hspace{3cm}}
	\vspace{0.6cm}
	
	% --- INSTRUCCIONES ---
	\noindent \textit{Instrucciones: Resuelve este taller utilizando lo que aprendiste en el TAU 3. Analiza cada pregunta con cuidado antes de responder.}
	\vspace{0.5cm}
	
	% --- PREGUNTA 1 --- 
	\noindent \textbf{1. Los grandes sabios de la antigüedad:} En este TAU conocimos a dos antiguos astrónomos griegos que usaron métodos muy sencillos para calcular tamaños astronómicos. Relaciona cada sabio con su gran descubrimiento uniendo con una línea:
	
	\vspace{0.4cm}
	\begin{itemize}[itemsep=0.8cm, label={}]
		\item \textbf{Aristarco de Samos} \hfill $\circ$ Usó la sombra de unos palos clavados en el suelo al mediodía para calcular el tamaño exacto de la \textbf{Tierra}.
		\item \textbf{Eratóstenes de Cirene} \hfill $\circ$ Usó la sombra que la Tierra proyecta durante un eclipse para calcular el tamaño de la \textbf{Luna}.
	\end{itemize}
	
	\solucion{\textbf{Aristarco de Samos} $\rightarrow$ Usó la sombra de la Tierra durante un eclipse para medir la Luna.\\ \textbf{Eratóstenes de Cirene} $\rightarrow$ Usó la sombra de los palos al mediodía para medir la Tierra.}
	\vspace{0.5cm}
	
	% --- PREGUNTA 2 --- 
	\noindent \textbf{2. El método de Aristarco:} ¿Cómo se dio cuenta Aristarco de Samos de que la Luna era más pequeña que la Tierra?
	\vspace{0.3cm} 
	\begin{itemize}[itemsep=0.4cm, label={}] 
		\item[ $\square$ ] Porque viajó hasta la Luna y la midió con una regla gigante.
		\item[ $\square$ ] Porque observó un eclipse de Luna y notó que la curva o ``barriga'' de la sombra de la Tierra era mucho más ancha que la propia Luna.
		\item[ $\square$ ] Porque el Sol iluminó la Luna por la mitad y se vio que era pequeña.
	\end{itemize}
	
	\solucion{La respuesta correcta es la segunda opción: Observó la ``barriga'' de la sombra de la Tierra.}
	\vspace{0.5cm}
	
	% --- PREGUNTA 3 --- 
	\noindent \textbf{3. Tamaños reales vs. Tamaños aparentes:} Lee cada afirmación y marca con una \textbf{X} si es Verdadera (V) o Falsa (F):
	
	\vspace{0.3cm} 
	\begin{itemize}[itemsep=0.6cm, label={}] 
		\item \textbf{A.} El ``tamaño real'' de un objeto cambia constantemente si te acercas o te alejas de él. \hfill  $\square$  V \quad  $\square$  F 
		
		\item \textbf{B.} Según la Regla de los Tamaños Aparentes, si te alejas de un objeto a una distancia 10 veces mayor, el objeto se verá 10 veces más pequeño. \hfill  $\square$  V \quad  $\square$  F 
	\end{itemize}
	
	\solucion{\textbf{A es Falsa} (el tamaño real jamás cambia, lo que cambia es el tamaño aparente).\\ \textbf{B es Verdadera} (Esa es la definición exacta de la regla vista en el libro).}
	\vspace{0.5cm}
	
	% --- PREGUNTA 4 --- 
	\noindent \textbf{4. El cálculo de Eratóstenes:} Eratóstenes descubrió que la distancia entre las antiguas ciudades de Siena y Alejandría era de \textbf{800 kilómetros}, y que ese pedazo correspondía exactamente a \textbf{1/50} (una cincuentava parte) del perímetro o circunferencia total de la Tierra. 
	
	\vspace{0.2cm}
	\noindent Con base en este dato, calcula cuál fue el resultado final que obtuvo Eratóstenes para la circunferencia de la Tierra. 
	\vspace{0.1cm}
	
	\noindent \textit{(Muestra tu operación matemática en el recuadro):}
	\vspace{0.2cm}
	
	\begin{tcolorbox}[colframe=black, colback=white, height=3.5cm, width=0.8\textwidth]
		\ifthenelse{\boolean{mostrarRespuestas}}{
			\textcolor{red!70!black}{
				\textbf{Operación esperada:}\\
				$800 \text{ kilómetros} \times 50 = 40.000 \text{ kilómetros}$\\
				\vspace{0.2cm}\\
				\textbf{Resultado:} Eratóstenes calculó que la circunferencia de la Tierra medía 40.000 km.}
		}{}
	\end{tcolorbox}
	\vspace{0.4cm}
	
	% --- PREGUNTA 5 --- 
	\noindent \textbf{5. ¿A qué distancia está la Luna?:} Al final del TAU se nos explica que la distancia entre la Tierra y la Luna es de aproximadamente \textbf{384.000 kilómetros}. Para hacer el cálculo mental más fácil, el Autor propone redondear esa cifra a \textbf{400.000 kilómetros}. 
	
	\vspace{0.2cm}
	\noindent Sabiendo que darle una vuelta completa a la Tierra (su circunferencia) son \textbf{40.000 kilómetros}, ¿a cuántas vueltas a la Tierra equivale un viaje a la Luna?
	\vspace{0.1cm}
	
	\noindent \textit{(Muestra tu operación matemática y escribe la respuesta final):}
	\vspace{0.2cm}
	
	\begin{tcolorbox}[colframe=black, colback=white, height=3.5cm, width=0.8\textwidth]
		\ifthenelse{\boolean{mostrarRespuestas}}{
			\textcolor{red!70!black}{
				\textbf{Operación esperada:}\\
				$400.000 \text{ km} \div 40.000 \text{ km} = 10 \text{ vueltas}$\\
				\vspace{0.4cm}\\
				\textbf{Respuesta final:} El viaje a la Luna equivale a \textbf{10} vueltas a la Tierra.}
		}{
			\vspace*{1.5cm}
			\textbf{Respuesta:} El viaje a la Luna equivale a \underline{\hspace{3cm}} vueltas a la Tierra.
		}
	\end{tcolorbox}
	
\end{document}